\loai{Biết cụ thể các thông số, tìm thông số còn lại}
\begin{ex} %[1P5B5-2][Loai1]
	Ở thời kì nén của một động cơ đốt trong 4 kì, nhiệt độ của hỗn hợp khí tăng từ 47\doC đến 367\doC, còn thể tích của khí giảm từ 1,8 lít đến 0,3 lít. Áp suất của khí lúc bắt đầu nén là 100 kPa. Coi hỗn hợp khí như chất khí thuần nhất, áp suất cuối thời kì nén là
	\choice
	{$1,5.10^6\ Pa$}
	{\True $1,2.10^6\ Pa$}
	{$1,8.10^6\ Pa$}
	{$2,4.10^6\ Pa$}
	\loigiai{Áp dụng phương trình trạng thái
		\begin{align*}
			\dfrac{p_1V_1}{T_1}=\dfrac{p_2V_2}{T_2}\Leftrightarrow \dfrac{100.10^3.1,8}{47+273}=\dfrac{p_2.0,3}{367+273}\Rightarrow p_2=1,2.10^6\ Pa.
	\end{align*}}
\end{ex}
\begin{ex} %[1P5B5-2][Loai1]
	Trong xilanh của một động cơ đốt trong có $2\ dm^3$ hỗn hợp khí dưới áp suất 1 atm và nhiệt độ 47\doC. Pít tông nén xuống làm cho thể tích của hỗn hợp khí chỉ còn $0,2\ dm^3$ và áp suất tăng lên tới 15 atm. Tính nhiệt độ của hỗn hợp khí nén.
	\choice
	{$107\ ^{\circ}C$}
	{$173\ ^{\circ}C$}
	{\True $207\ ^{\circ}C$}
	{$273\ ^{\circ}C$}
	\loigiai{Áp dụng phương trình trạng thái:
		\begin{align*}
			\dfrac{p_1V_1}{T_1}=\dfrac{p_2.V_2}{T_2}\Rightarrow T_2=480K\Rightarrow t_2=207\doC.
	\end{align*}}
\end{ex}
\begin{ex} %[1P5B5-2][Loai1]
	Trong phòng thí nghiệm, người ta điều chế được $60\ cm^3$ khí hidro ở áp suất 750 mmHg và nhiệt độ $27\doC$. Tính thể tích của lượng khí trên ở điều kiện chuẩn (áp suất 760 mmHg và nhiệt độ $0\doC$).
	\choice
	{$50\ cm^3$}
	{$59,21\ cm^3$}
	{\True $53,88\ cm^3$}
	{$60,5\ cm^3$}
	\loigiai{\begin{itemize}
			\item Trạng thái 1: $p_1 = 750\ mmHg,\ {V_1} = 60\ cm^3,\ T_1 = 300\ K$.
			\item Trạng thái 2: $p_2 = 760\ mmHg,\ V_2,\ T_1 = 273\ K$.
			\item Phương trình trạng thái: $\dfrac{750.60}{300} = \dfrac{760.{V_2}}{273} \Rightarrow V_2 = 53,88\ cm^3$.
	\end{itemize}}
\end{ex}
\begin{ex} %[1P5B5-2][Loai1]
	Nén 10 lít khí ở nhiệt độ 27\doC để thể tích của nó giảm chỉ còn 4 lít, quá trình nén nhanh nên nhiệt độ tăng đến 60\doC. Áp suất khí đã tăng bao nhiêu lần?
	\choice
	{\True 2,78}
	{3,2}
	{2,24}
	{2,85}
	\loigiai{Áp dụng phương trình trạng thái
		\begin{align*}
			\dfrac{p_1V_1}{T_1}=\dfrac{p_2V_2}{T_2}\Leftrightarrow \dfrac{p_2}{p_1}=\dfrac{V_1T_2}{V_2T_1}=\dfrac{10.(60+273)}{4.(27+273)}=2,78.
	\end{align*}}
\end{ex}
\begin{ex} %[1P5K5-2][Loai1]
	Trong xilanh của một động cơ có chứa một lượng khí ở nhiệt độ 47\doC và áp suất 0,7 atm.
	\begin{itemize}
		\item[a)] Sau khi bị nén thể tích của khí giảm đi 5 lần và áp suất tăng lên tới 8 atm. Tính nhiệt độ của khí ở cuối quá trình nén?
		\item[b)] Người ta tăng nhiệt độ của khí lên đến 273\doC và giữ pit-tông cố định thì áp suất của khí khi đó là bao nhiêu?
	\end{itemize}
	\loigiai{a) Áp dụng PTTT khí lý tưởng, ta có:
		\begin{align*}
			\dfrac{p_1V_1}{T_1}=\dfrac{p_2V_2}{T_2}\Rightarrow T_2=\dfrac{8V_1.320}{5.0,7V_1}=731\ K.
		\end{align*}
		b) Vì pít- tông được giữ cố định nên đó là quá trình đẳng tích.\\
		Theo định luật Sác - lơ, ta có:
		\begin{align*}
			\dfrac{p_1}{T_1}=\dfrac{P_3}{T_3}\Rightarrow p_3=\dfrac{p_1.T_3}{T_1}=\dfrac{546.0,7}{320}=1,19\ atm.
	\end{align*}}
\end{ex}
\loai{Các thông số thay đổi một lượng so với trước}
\begin{ex} %[1P5K5-2][Loai2]
	Trong một động cơ điezen, khối khí có nhiệt độ ban đầu là 32\doC được nén để thể tích giảm bằng $\dfrac{1}{16}$ thể tích ban đầu và áp suất tăng bằng 48,5 lần áp suất ban đầu. Nhiệt độ khối khí sau khi nén sẽ bằng
	\choice
	{97\doC}
	{\True 652\doC}
	{1552\doC}
	{132\doC}
	\loigiai{Áp dụng phương trình trạng thái khí lý tưởng
		\begin{align*}
			\dfrac{p_1V_1}{T_1}=\dfrac{p_2V_2}{T_2}\Leftrightarrow \dfrac{p_1V_1}{32+273}=\dfrac{48,5p_1.V_1/16}{t_2+273}\Rightarrow t_2=652\doC.
	\end{align*}}
	%<MyLT>
\end{ex}
\begin{ex} %[1P5B5-2][Loai2]
	Một khối khí lí tưởng đang ở áp suất 2\ atm thì được nung nóng đến khi nhiệt độ tuyệt đối tăng lên 4 lần và thể tích tăng lên 2 lần. Áp suất của khối khí sau khi nung là
	\choice
	{0,25\ atm}
	{\True 4\ atm}
	{16\ atm}
	{2\ atm}
	\loigiai{Phương trình trạng thái $\dfrac{pV}{T}= const\Rightarrow$ T tăng 4 lần, V tăng 2 lần thì p giảm 2 lần $\Rightarrow p' = 1\ atm$}
\end{ex}
\begin{ex} %[1P5B5-2][Loai2]
	Thể tích của hỗn hợp khí trong xilanh là $2\ dm^3$, nhiệt độ là $47\doC$ và áp suất ban đầu là 1\ atm. Tính nhiệt độ của hỗn hợp khí trong xilanh khi pít-tông nén khí trong xilanh làm thể tích giảm đi 10 lần, áp suất tăng đến 15 atm.
	\choice
	{423\ K}
	{\True 480\ K}
	{517\ K}
	{470\ K}
	\loigiai{\begin{itemize}
			\item Trạng thái 1: $T_1 = 320\ K,\ p_1 = 1\ atm,\ V_1 = 2\ dm^3$.
			\item Trạng thái 2: $T_2,p_2= 15\ atm,\ {V_2} = 0,2\ dm^3$.
			\item Phương trình trạng thái: $\dfrac{p_1V_1}{T_1} = \dfrac{p_2V_2}{T_2} \Rightarrow \dfrac{1.2}{320} = \dfrac{15.0,2}{T_2}\Rightarrow T_2 = 480\ K$.
	\end{itemize}}
\end{ex}
\begin{ex} %[1P5K5-2][Loai2]
	Nếu thể tích của một lượng khí giảm đi $\dfrac{1}{10}$, áp suất tăng $\dfrac{1}{5}$ và nhiệt độ tăng thêm $16\doC$ so với ban đầu. Tính nhiệt độ ban đầu của khí.
	\choice
	{100 K}
	{\True 200 K}
	{150 K}
	{250 K}
	\loigiai{Từ phương trình trạng thái khí lý tưởng:
		\begin{align*}
			\dfrac{p_1V_1}{T_1}=\dfrac{p_2.V_2}{T_2}\Rightarrow T_1=200\ K.
	\end{align*}}
\end{ex}
\loai{Bài toán liên quan khối lượng riêng của lượng khí}
\begin{ex} %[1P5K5-2][Loai3]
	Tính khối lượng riêng của không khí ở $100\doC$ và áp suất $2.10^5\ Pa$. Biết khối lượng riêng của không khí ở $0\doC$ là $1,29\ kg/m^3$.
	\choice
	{$1,53\ kg/m^3$}
	{$1,79\ kg/m^3$}
	{$2\ kg/m^3$}
	{\True $1,86\ kg/m^3$}
	\loigiai{\begin{itemize}
			\item Trạng thái 1: $p_1 = 2.10^5\ Pa$, $V_1 = \dfrac{m}{D_1},\ T_1 = 373\ K$.
			\item Trạng thái 2: $p_2 = 1,013.10^5\ Pa,\ V_2= m/1,29,\ T_2 = 273\ K$.
			\item Phương trình trạng thái: $\dfrac{2.10^5.\dfrac{m}{D_1}}{373} = \dfrac{1,013.10^5.\dfrac{m}{1,29}}{273} \Rightarrow D_1 = 1,86\ kg/m^3$.
	\end{itemize}}
\end{ex}
\begin{ex} %[1P5K5-2][Loai3]
	Không khí ở áp suất $10^5\ Pa$, nhiệt độ $0\doC$ có khối lượng riêng $1,19\ kg/m^3$. Tính khối lượng riêng của không khí ở áp suất $2.10^5\ Pa$, nhiệt độ $100\doC$.
	\choice
	{\True $1,74\ kg/m^3$}
	{$2,38\ kg/m^3$}
	{$1,89\ kg/m^3$}
	{$1,29\ kg/m^3$}
	\loigiai{\begin{itemize}
			\item Trạng thái 1: $T_1 = 273\ K,\ {p_1} = {10^5}\ Pa,\ {V_1} = \dfrac{m}{D_1}$.
			\item Trạng thái 2: $T_2 = 373\ K,\ {p_2} = {2.10^5}\ Pa,\ {V_2} = \dfrac{m}{D_2}$.
			\item Phương trình trạng thái: $\dfrac{p_1V_1}{T_1} = \dfrac{p_2V_2}{T_2} \Rightarrow \dfrac{p_1}{T_1D_1} = \dfrac{p_2}{T_2D_2}\Rightarrow {D_2} = \dfrac{T_1D_1p_2}{p_1T_2}= \dfrac{273.1,19.2.10^5}{10^5.373} = 1,74\ kg/m^3$.
	\end{itemize}}
\end{ex}
\loai{Bài toán bơm nén khí vào bình}
\begin{ex} %[1P5B5-2][Loai4]
	Một máy nén khí ở áp suất 1\ atm mỗi lần nén được 4 lít khí ở nhiệt độ $27\doC$ vào trong bình chứa thể tích $2\ m^3$ áp suất ban đầu 1\ atm. Tính áp suất khí bên trong bình chứa sau 1000 lần nén khí biết nhiệt độ trong bình sau 1000 lần nén là $42\doC$.
	\choice
	{3,3\ atm}
	{\True 3,15\ atm}
	{1,5\ atm}
	{2,7\ atm}
	\loigiai{\begin{itemize}
			\item Trạng thái 1: $T_1 = 300\ K,\ p_1 = 1\ atm,\ V_1 = 2 + 4.1000.10^{-3} = 6\ m^3$.
			\item Trạng thái 2: $T_2 = 315\ K,\ p_2,\ {V_2} = 2\ m^3$.
			\item Phương trình trạng thái: $\dfrac{p_1V_1}{T_1} = \dfrac{p_2V_2}{T_2} \Rightarrow \dfrac{1.6}{300} = \dfrac{{{p_2}.2}}{315}\Rightarrow p_2 = 3,15\ atm$.
	\end{itemize}}
\end{ex}
\loai{Bài toán dùng bình khí nén bơm bóng}
\begin{ex} %[1P5G5-2][Loai5]
	Một bình kín dung tích không đổi 50 lít chứa khí Hyđrô ở áp suất 5 MPa và nhiệt độ 37\doC, dùng bình này để bơm bóng bay, mỗi quả bóng bay được bơm đến áp suất $1,05.10^5\ Pa$, dung tích mỗi quả là 10 lít, nhiệt độ khí nén trong bóng là 12\doC. Bình đó bơm được bao nhiêu quả bóng bay?
	\choice
	{200}
	{150}
	{\True 214}
	{188}
	\loigiai{\begin{itemize}
			\item Trạng thái 1 có: $\begin{cases}
				V_1 = 50\ (lit) \\
				p_1 = 5.10^6\ (Pa) \\
				T_1=37+273=310\ (K)
			\end{cases}$
			\item Trạng thái 2 có: $\begin{cases}
				V_2 = ?\ (lit) \\
				p_2 = 1,05.10^5\ (Pa) \\
				T_2=12+273=285\ (K)
			\end{cases}$\\
			\item Ở nhiệt độ 12\doC và áp suất $1,05.10^5$ Pa thì thể tích khí của lượng Hidro chứa trong chai là:
			\begin{align*}
				\dfrac{p_1V_1}{T_1}=\dfrac{p_2V_2}{T_2}\Leftrightarrow V_2=\dfrac{p_1V_1T_2}{T_1p_2}=\dfrac{5.10^6.50.285}{310.1,05.10^5}=2189\ \text{lít}.
			\end{align*}
			\item Vì không thể lấy hết lượng khí ở trong bình nên ta phải trừ bớt 50 lít, khi đó số bóng bơm được là
			\begin{align*}
				\dfrac{2189-50}{10}\approx 214\quad \text{quả}.
			\end{align*}
	\end{itemize}}
\end{ex}
\loai{Xác định các thông số từ đồ thị các đẳng quá trình}
\begin{ex} %[1P5K5-2][Loai6]
	immini[thm]{Một khối khí thực hiện một chu trình như hình vẽ. Cho $T_2 = 900\ K$, tính $V_2$ và $T_3$.
		\choice
		{$V_2 = 2$ lít, $T_3 = 300\ K$}
		{\True $V_2 = 6$ lít, $T_3 = 300\ K$}
		{$V_2 = 6$ lít, $T_3 = 900\ K$}
		{$V_2 = 2$ lít, $T_3 = 900\ K$}}{\begin{tikzpicture}[scale=1,thick,>=stealth']
			\draw[->] (0,0)--(4,0)node[below]{$V\ (\ell)$};
			\node at (o) [below left]{\normalsize{O}};
			\draw[->] (0,0)--(0,3.5)node[right]{$p\ (Pa)$};
			\draw[dashed](1,0)node[below]{$2$}--(1,3)
			(0,1)node[left]{$2.10^5$}--(3,1)
			(0,3)node[left]{$6.10^5$}--(1,3)(3,1)--(3,0);
			\draw[undirected,blue, ultra thick](1,1)node[above left,inner sep=1pt,black]{\small{(3)}}--(3,1)node[right,black]{\small{(2)}};
			\draw[undirected,blue, ultra thick](3,1) to [bend left=40] (1,3)node[right,black]{\small{(1)}};
			\draw[directed,blue, ultra thick](1,1) -- (1,3);
	\end{tikzpicture}}
	\loigiai{\begin{itemize}
			\item Từ $(1)\to(2)$ là quá trình đẳng nhiệt $\Rightarrow {p_1}{V_1} = {p_2}{V_2} \Leftrightarrow {6.10^5}.2 = {2.10^5}.{V_2} \Rightarrow V_2 = 6$ lít.
			\item Từ $(2)\to(3)$ là quá trình đẳng áp $\Rightarrow \dfrac{V_2}{T_2} = \dfrac{V_3}{T_3} \Leftrightarrow \dfrac{6}{900} = \dfrac{2}{T_3} \Rightarrow T_3 = 300\ K$.
	\end{itemize}}
\end{ex}
\begin{ex} %[1P5K5-2][Loai6]
	immini[thm]{Một khí lí tưởng ban đầu ở trạng thái 1 có thể tích 20 lít, áp suất 1\ atm, nhiệt độ 300\ K biến đổi trạng trạng thái qua hai quá trình liên tiếp được biểu diễn bằng đồ thị trong hệ trục (p, V) như hình vẽ bên. Tính nhiệt độ của khối khí ở cuối quá trình (2) và (3).
		\choice
		{$T_2 = 200\ K,\ T_3 = 900\ K$}
		{\True $T_2 = 450\ K,\ T_3 = 900\ K$}
		{$T_2 = 300\ K,\ T_3 = 900\ K$}
		{$T_2 = 450\ K,\ T_3 = 600\ K$}}{\begin{tikzpicture}[scale=1,thick,>=stealth']
			\node at (o) [below left]{\normalsize{O}};
			\draw[->] (0,0)--(4,0)node[below]{$V\ (\ell)$};
			\draw[->] (0,0)--(0,3.5)node[right]{$p\ (atm)$};
			\draw[dashed](0,1)node[left]{$1$}--(2,1)node[right,inner sep=1pt]{\small(2)}--(2,0)node[below]{$30$};
			\draw[dashed](0,3)node[left]{$2$}--(2,3)node[right,inner sep=1pt]{\small(3)};
			\draw[dashed](1.3,0)node[below]{$20$}--(1.3,1);
			\draw[arrow data={0.5}{stealth},blue, ultra thick] (1.3,1)node[below left,inner sep=1pt,black]{\small(1)}--(2,1);
			\draw[arrow data={0.5}{stealth},blue, ultra thick] (2,1)--(2,3);
	\end{tikzpicture}}
	\loigiai{\begin{itemize}
			\item Qúa trình từ $(1)\to(2)$ là quá trình đẳng áp $\Rightarrow \dfrac{V_1}{T_1} = \dfrac{V_2}{T_2} \Leftrightarrow \dfrac{20}{300} = \dfrac{30}{T_2} \Rightarrow T_2 = 450\ K$.
			\item Qúa trình từ $(2)\to(3)$ là quá trình đẳng tích $\Rightarrow \dfrac{p_2}{T_2} = \dfrac{p_3}{T_3} \Leftrightarrow \dfrac{1}{450} = \dfrac{2}{T_3} \Rightarrow T_3 = 900\ K$.
	\end{itemize}}
\end{ex}
\begin{ex} %[1P5K5-2][Loai6]
	immini[thm]{Sự biến đổi trạng thái của một khối khí lí tưởng được mô tả như hình vẽ, $V_1 = 3$ lít, $V_3 = 6$ lít. Tính $p_2,\ T_1,\ V_2$.
		\choice
		{$p_2 = 1\ atm,\ T_1 = 300\ K,\ V_2 = 1$ lít}
		{$p_2 = 2\ atm,\ T_1 = 600\ K,\ V_2 = 3$ lít}
		{\True $p_2 = 2\ atm,\ T_1 = 300\ K,\ V_2 = 3$ lít}
		{$p_2 = 1\ atm,\ T_1 = 600\ K,\ V_2 = 1$ lít}}{\begin{tikzpicture}[scale=0.8,thick,>=stealth']
			\tkzDefPoints{0/0/o, 4/3/a, 4/1.5/b, 2/1.5/c}
			\node at (o) [below left]{\normalsize{O}};
			\draw[->] (0,0)--(5,0)node[below]{$T$};
			\draw[->] (0,0)--(0,3.5)node[right]{$p\ (atm)$};
			\draw[dashed](0,1.5)--(c)--(2,0)(0,3)--(a)(b)--(4,0)(o)--(a);
			\tkzDrawSegments[blue,very thick,arrow data={0.5}{stealth}](a,b b,c c,a)
			\node at (a)[above right]{\small(2)};
			\node at (b)[right]{\small(3)};
			\node at (c)[above left]{\small (1)};
			\node at (0,3)[left]{\small $p_{2}$};
			\node at (0,1.5)[left]{\small $p_{1}$};
			\node at (2,0)[below]{\small $T_{1}$};
			\node at (4,0)[below]{\small $600$};
	\end{tikzpicture}}
	\loigiai{\begin{itemize}
			\item Từ $(1) \to (2)$: đẳng tích $\Rightarrow V_2 = V_1 = 3$ lít.
			\item Từ $(2) \to (3)$: đẳng nhiệt $\Rightarrow {T_2} = {T_3} = 600\ K$.
			\item Từ $(3) \to (1)$: đẳng áp $\Rightarrow {p_1} = {p_3} = 1\ atm$.
			\item Qúa trình từ $(3)$ đến $(1)$ là đẳng áp $\Rightarrow \dfrac{V_3}{T_3} = \dfrac{V_1}{T_1} \Leftrightarrow \dfrac{6}{600} = \dfrac{3}{T_1} \Rightarrow T_1 = 300\ K$.
			\item Qúa trình từ $(2)$ đến $(3)$ là đẳng nhiệt $\Rightarrow {p_2}{V_2} = {p_3}V{_3} \Leftrightarrow {p_2}.3 = 1.6 \Rightarrow p_2 = 2\ atm$.
	\end{itemize}}
\end{ex}
\begin{ex} %[1P5K5-2][Loai6]
	immini[thm]{Trên hình vẽ biểu diễn sự biến đổi trạng thái của một lượng khí lí tưởng trong hệ tọa độ (p,V). Cho biết $t_1 = 27\doC$. Tính nhiệt độ cuối $T_3$ của lượng khí đó.
		\choice
		{\True 900 K}
		{$427\doC$}
		{$276\doC$}
		{372 K}}{\begin{tikzpicture}[scale=1,thick,>=stealth']
			\node at (o) [below left]{\normalsize{O}};
			\draw[->] (0,0)--(4,0)node[below]{$V\ (\ell)$};
			\draw[->] (0,0)--(0,3.5)node[right]{$p\ (atm)$};
			\draw[dashed](0,1)node[left]{$1$}--(1.3,1)--(1.3,0)node[below]{$10$};
			\draw[dashed](0,2)node[left]{$2$}--(2,2)node[right,inner sep=1pt]{\small(3)}--(2,0)node[below]{$15$};
			\draw[arrow data={0.5}{stealth},blue, ultra thick] (1.3,1)node[below left,inner sep=1pt,black]{\small(1)}--(1.3,2)node[above,inner sep=1pt,black]{\small(2)};
			\draw[arrow data={0.5}{stealth},blue, ultra thick] (1.3,2)--(2,2);
	\end{tikzpicture}}.
	\loigiai{\begin{align*}
			\begin{cases}
				{p_1} = 1\,\ atm\\
				{V_1} = 10\ \ell \\
				{T_1} = 200\ K
			\end{cases}
			\to
			\begin{cases}
				{p_2} = 2\ atm\\ {V_2} = 10\ \ell \\
				{T_2}
			\end{cases}
			\to
			\begin{cases}
				{p_3} = 2\,\ atm\\
				{V_3} = 15\ \ell \\
				{T_3}
			\end{cases}
		\end{align*}
		\begin{itemize}
			\item Quá trình từ $(1)$ đến $(2)$ là quá trình đẳng tích $\Rightarrow \dfrac{1}{300} = \dfrac{2}{T_2} \Rightarrow T_2 = 600\ K$.
			\item Quá trình từ $(2)$ đến $(3)$ là quá trình đẳng áp $\Rightarrow \dfrac{10}{600} = \dfrac{15}{T_3} \Rightarrow T_3 = 900\ K$.
	\end{itemize}}
\end{ex}
\loai{Xác định thông số từ đồ thị không phải đẳng quá trình}
\begin{ex} %[1P5G5-2][Loai7]
	immini[thm]{Một lượng khí biến đổi theo chu trình biểu diễn bởi đồ thị. Cho biết ${p_1} = {p_3},\ {V_1} = 1\ m^3,\ {V_2} = 4\ m^3,\ T_1 = 100\ K,\ T_4 = 300\ K$. Tìm $V_3$.
		\choice
		{$2\ m^3$}
		{$1,6\ m^3$}
		{$2,5\ m^3$}
		{\True $2,2\ m^3$}}
	{\begin{tikzpicture}[scale=0.8,thick,>=stealth']
			\tkzDefPoints{0/0/o, 1/1/a, 1/4/b, 2.2/2.2/c, 3/1/d}
			\node at (o) [below left]{\normalsize{O}};
			\draw[->] (0,0)--(4,0)node[below]{$T$};
			\draw[->] (0,0)--(0,5)node[left]{$V$};
			\draw[dashed](0,1)--(a)--(1,0)(0,4)--(b)(o)--(c)(d)--(3,0);
			\tkzDrawSegments[blue,very thick,arrow data={0.5}{stealth}](a,b b,d d,a)
			\node at (a)[above left]{\small(1)};
			\node at (b)[above]{\small(2)};
			\node at (c)[above right]{\small (3)};
			\node at (d)[right]{\small (4)};
			\node at (0,4)[left]{\small $V_{2}$};
			\node at (0,1)[left]{\small $V_{1}$};
			\node at (1,0)[below]{\small $T_{1}$};
			\node at (3,0)[below]{\small $T_{4}$};
	\end{tikzpicture}}
	\loigiai{\begin{itemize}
			\item Qúa trình $(1)-(2)$: đẳng nhiệt, ${T_2} = {T_1} = 100\ K,\ {V_2} = 4\ m^3$.
			\item Qúa trình $(4)-(1)$: đẳng tích, ${V_4} = {V_1} = 1\ m^3,\ {T_4} = 300\ K$.
			\item Qúa trình từ $(2)-(4)$: $V = aT + b$
			\begin{itemize}
				\item Trạng thái $(2)$: $4 = 100a + b-$
				\item Trạng thái $(4)$: $1 = 300a +b\Rightarrow a = -0,015, b = 5,5\Rightarrow V = - 0,015T + 5,5$
			\end{itemize}
			\item Qúa trình từ $(1)-(3)$ là đẳng áp $\Rightarrow V = \dfrac{V_1}{T_1}T = \dfrac{1}{100}T = 0,01T$
			\item Vì trạng thái $(3)$ là giao điểm của hai đường $(2)-(4)$ và $(1)-(3)$ nên:
			\begin{align*}
				-0,015T_3 + 5,5 = 0,01T_3\Rightarrow {T_3} = 220\ K\Rightarrow V3=0,01.220=2,2m3{V_3} = 0,01.220 = 2,2\ m^3.
			\end{align*}
	\end{itemize}}
\end{ex}
\loai{Bài toán pittông hoặc giọt thuỷ ngân dịch chuyển}
\begin{ex} %[1P5G5-2][Loai8]
	immini[thm]{Bình kín được ngăn làm hai phần bằng nhau (A và B) bằng tấm cách nhiệt có thể dịch chuyển được. Biết mỗi bên có chiều dài 30 cm và nhiệt độ của khí trong bình là $27\doC$, xác định khoảng dịch chuyển của tấm cách nhiệt khi nung nóng phần A thêm $10\doC$, đồng thời hạ nhiệt độ phần B đi $10\doC$.
		\choice
		{3 cm}
		{4 cm}
		{\True 1 cm}
		{2 cm}}{\begin{tikzpicture}[scale=1,thick,>=stealth']
			\draw(0,0)--++(0:4)--++(-90:1)--++(180:4)--cycle;
			\draw[pattern=north east lines](1.9,0)--++(0:0.2)--++(-90:1)--++(180:0.2)--cycle;
			\node at (1,-0.5){$A$};
			\node at (3,-0.5){$B$};
	\end{tikzpicture}}
	\loigiai{\begin{itemize}
			\item Gọi h là chiều cao của bình, y là chiều rộng của bình, x là khoảng vách ngăn dịch chuyển
			\item Phần A:
			\begin{itemize}
				\item Trạng thái 1: ${V_0} = h{\ell _0}y,\ p_0,\ T_0 = 300\ K$.
				\item Trạng thái 2: ${V_{2A}} = h\left( {{\ell _0} + x} \right)y,\ {p_{2A}},\ {T_{2A}} = 310\ K$.
				\item Theo phương trình trạng thái: $\dfrac{{{p_0}.h{\ell _0}y}}{300} = \dfrac{{p_{2A}.h\left( {{\ell _0} + x} \right)y}}{310}\quad (1)$.
			\end{itemize}
			\item Phần B:
			\begin{itemize}
				\item Trạng thái 1: $V_0 = h{\ell _0}y,\ p_0,\ T_0 = 300\ K$.
				\item Trạng thái 2: $V_{2B} = h\left( {{\ell _0} - x} \right)y,\ {p_{2B}},\ {T_{2B}} = 290\ K$.
				\item Theo phương trình trạng thái: $\dfrac{{{p_0}.h{\ell _0}y}}{300} = \dfrac{{p_{2B}.h\left( {{\ell _0} - x} \right)y}}{290}\quad (2).$
			\end{itemize}
			\item Để vách ngăn nằm cân bằng sau khi nung nóng một bên và làm lạnh một bên thì áp suất của phần A và phần B sau khi nung nóng phải bằng nhau $\Rightarrow {p_{2A}} = {p_{2B}}\quad (3)$
			\item Từ $(1)$, $(2)$ và $(3) \Rightarrow x = 1\ cm$.
	\end{itemize}}
\end{ex}
\begin{ex} %[1P5G5-2][Loai8]
	Một xilanh kín được chia làm 2 phần, mỗi phần dài 52 cm và ngăn cách nhau bằng pít-tông cách nhiệt. Mỗi phần chứa một lượng khí giống nhau ở $27\doC$, 750 mmHg. Khi nung nóng một phần thêm $50\doC$ thì pít-tông dịch chuyển một đoạn bằng 4 cm. Tìm áp suất sau khi nung.
	\choice
	{696,4 mmHg}
	{\True 812,5 mmHg}
	{758,3 mmHg}
	{524,8 mmHg}
	\loigiai{\begin{itemize}
			\item Ban đầu, khí trong mỗi phần xilanh có thể tích $V = 52S$, áp suất p = 750\ mmHg, nhiệt độ $T = 300\ K$.
			\item Sau khi nung:Phần khí bị nung: ${V_1} = 56S,\ p_1,\ T_1 = 350\ K$.
			\item Phần khí không bị nung: ${V_2} = 48S,\ {p_2} = {p_1},\ T_2 = 300\ K$.
			\item Vì nhiệt độ của lượng khí trong phần xilanh không bị nung nóng khối đổi, đều bằng 300\ K, nên áp dụng định luật Bôi-lơ Ma-ri-ốt cho lượng khí này, ta được:
			\begin{align*}
				pV = {p_2}{V_2} \Rightarrow 750.52S = {p_2}.48S \Rightarrow p_2 = 812,5\ mmHg.
			\end{align*}
	\end{itemize}}
\end{ex}
\begin{ex} %[1P5G5-2][Loai8]
	Hai bình chứa cùng một lượng khí nối với nhau bằng một ống nằm ngang tiết diện $0,4\ cm^2$, ngăn cách nhau bằng một giọt thủy ngân trong suốt. Ban đầu mỗi phần có nhiệt độ $27\doC$, thể tích 0,3 lít. Tính khoảng dịch chuyển của một giọt thủy ngân khi nhiệt độ bình I tăng thêm $2\doC$, bình II giảm $2\doC$. Coi bình dãn nở không đáng kể.
	\choice
	{4 cm}
	{\True 5 cm}
	{3 cm}
	{2 cm}
	\loigiai{\begin{itemize}
			\item Ban đầu, khí trong mỗi phần xilanh có thể tích $V = 0,3\ \ell$, áp suất p, nhiệt độ $T = 300\ K$
			\item Khi nhiệt độ các bình thay đổi thì: Khí trong bình có: ${V_1} = V + Sx,\ p_1,\ T_1 = T + 2$.
			\item Khí trong bình II có: ${V_2} = V - Sx,\ {p_2} = {p_1},\ T_2 = T - 2$.
			\item Vì khí trong hai bình giống nhau, áp suất khí trong hai bình khí cân bằng nhau nên
			\begin{align*}
				\dfrac{V_2}{V_1} = \dfrac{T_2}{T_1} = \dfrac{302}{298}\quad (1)
			\end{align*}
			\item Lại có: ${V_1} + {V_2} = 2V = 0,6\quad (2).$
			\item Từ $(1)$ và $(2) \Rightarrow V_1 = 0,298\ \ell,\ V_2 = 0,302\ \ell$ $\Rightarrow x = \dfrac{{{V_1} - V}}{S} = \dfrac{0,302.10^{- 3} - 0,3.10^{- 3}}{0,4.10^{- 4}} = 0,05 m = 5\ cm$.
	\end{itemize}}
\end{ex}
\begin{ex} %[1P5G5-2][Loai8]
	Một xilanh kín được chia làm 2 phần, mỗi phần dài 52\ cm và ngăn cách nhau bằng pít-tông cách nhiệt. Mỗi phần chứa một lượng khí giống nhau ở $27\doC,\ 750\ mmHg$.\ Khi nung nóng một phần thêm $50\doC$ thì pít-tông dịch chuyển một đoạn bằng
	\choice
	{1 cm}
	{3 cm}
	{\True 4 cm}
	{2 cm}
	\loigiai{\begin{itemize}
			\item Gọi độ dịch chuyển của pít-tông là x.
			\item Ban đầu, khí trong mỗi phần xilanh có thể tích $V = 52S$, áp suất $p = 750\ mmHg$, nhiệt độ $T = 300\ K$.
			\item Sau khi nung:
			\begin{itemize}
				\item Phần khí bị nung: ${V_1} = S\left( {52 + x} \right),\ p_1,\ T_1 = 350\ K$.
				\item Phần khí không bị nung: ${V_2} = S\left( {52 - x} \right),\ {p_2} = {p_1},\ T_2 = 300\ K$.
			\end{itemize}
			\item Vì lượng khí trong mỗi phần xilanh giống nhau, $p_1 = p_2$ nên
			\begin{align*}
				\dfrac{V_1}{V_2} = \dfrac{T_1}{T_2} \Leftrightarrow \dfrac{52 + x}{52 - x} = \dfrac{350}{300} \Rightarrow x = 4\ cm.
			\end{align*}
	\end{itemize}}
\end{ex}
\begin{ex} %[1P5G5-2][Loai8]
	Một xilanh kín chia làm hai phần bằng nhau bởi một pít-tông cách nhiệt. Mỗi phần có chiều dài 50 cm chứa một lượng khí giống nhau ở 30\doC. Nung nóng một phần lên 20\doC, còn phần kia làm lạnh đi 20\doC thì pít-tông dịch chuyển một đoạn là
	\choice
	{4 cm}
	{2 cm}
	{\True 3,3 cm}
	{0,5 cm}
	\loigiai{\begin{itemize}
			\item Khi pit tông đứng yên (trước và sau khi di chuyển) nên áp suất của khí hai bên pít-tông là như nhau.
			\item Áp dụng phương trình trạng thái cho khí trong mỗi phần xilanh:
			\item Phần khí bị nung nóng :
			\begin{align*}
				\dfrac{p_0V_0}{T_0}=\dfrac{p_1V_1}{T_1}\quad (1).
			\end{align*}
			\item Phần khí bị làm lạnh :
			\begin{align*}
				\dfrac{p_0V_0}{T_0}=\dfrac{p_2V_2}{T_2}\quad (2).
			\end{align*}
			\item Từ phương trình $(1),(2)$ và $p_1=p_2\Rightarrow \dfrac{V_1}{T_1}=\dfrac{V_2}{T_2}$.
			\item Gọi x là khoảng pit-tông dịch chuyển ta có
			\begin{align*}
				\dfrac{(\ell_0+x)S}{T_1}=\dfrac{(\ell_0-x)S}{T_2}\Rightarrow x=\dfrac{\ell_0(T_1-T_2)}{T_1+T_2}=\dfrac{50.(50-10)}{50+10}.
			\end{align*}
			\item Thay số ta được $x=3,3\ cm$.
	\end{itemize}}
\end{ex}
\begin{ex} %[1P5G5-2][Loai8]
	Một xilanh được chia làm hai phần, mỗi phần dài 42 cm và ngăn cách nhau bởi một pít-tông cách nhiệt. Mỗi phần xilanh chứa cùng một khối lượng khí giống nhau ở $27\doC$ và áp suất 1 atm. Cần phải nung nóng khí ở một phần của xilanh đến bao nhiêu để pít-tông có thể dịch chuyển 2 cm?
	\choice
	{\True 330 K}
	{$47\doC$}
	{$87\doC$}
	{303 K}
	\loigiai{\begin{itemize}
			\item Ban đầu, khí trong mỗi phần xilanh có thể tích $V = 42S$, áp suất $p = 1 at$, nhiệt độ $T = 300\ K$.
			\item Sau khi nung:Phần khí bị nung: ${V_1} = 44S,\ p_1,\ T_1$.
			\item Phần khí không bị nung: ${V_2} = 40S,\ {p_2} = {p_1},\ T_2 = 300\ K$.
			\item Vì lượng khí trong mỗi phần xilanh giống nhau, ${p_1} = {p_2}$ nên $\dfrac{V_1}{V_2} = \dfrac{T_1}{T_2} \Leftrightarrow \dfrac{44}{40} = \dfrac{T_1}{300} \Rightarrow T_1 = 330\ K$.
	\end{itemize}}
\end{ex}
\loai{Bài toán khác}
\begin{ex} %[1P5K5-2][Loai9]
	Một bóng thám không được chế tạo để có thể tăng bán kính lên tới 10 m khi bay ở tầng khí quyển có áp suất 0,03 atm và nhiệt độ 400 K. Bán kính của bóng khi bơm là bao nhiêu? Biết bóng được bơm khí ở áp suất 1 atm và nhiệt độ 600 K.
	\choice
	{\True 3,56 m}
	{5,6 m}
	{6,54 m}
	{5,36 m}
	\loigiai{\begin{itemize}
			\item Trạng thái 1: $p_1 = 0,03\ atm,\ V_1 = \dfrac{4}{3}.\pi .10^3\ m^3,\ T_1 = 400\ K$.
			\item Trạng thái 2: $p_1 = 1\ atm,\ V_2= \dfrac{4}{3}.\pi .r^3,\ T_2 = 600\ K$.
			\item Phương trình trạng thái: $\dfrac{{0,03.\dfrac{4}{3}.\pi {{.10}^3}}}{400} = \dfrac{{1.\dfrac{4}{3}\pi {r^3}}}{600} \Rightarrow r = 3,56\ m$.
	\end{itemize}}
\end{ex}


